\Needspace{12\baselineskip}
\originalchapter{Lab1：经验势的全部公开任务}
本章按3.320 Spring 2005的Lab1原题逐项组织，截止日期为2005年2月24日。输入包括手册明确印出的Au-LJ参数与FCC位置；所引用 \texttt{gulp1b.in} 的原Au-EAM函数/参数没有在本次取得文件中公开。本章保留EAM任务和输入需求，给出可复用方法，不用另一个势的结果冒充原势。新静态程序实际执行了LJ原给定系数的晶格、未弛豫空位及未弛豫表面子任务；它没有调用GULP。

\originalsection{晶格与总能量}
\begin{exercise}\label{task:lab1-1}
（原Lab1题1）(A)用 \texttt{gulp1a.in} 的LJ势计算FCC Au的平衡晶格常数和eV总能量：先用首行opti自动弛豫，再改single，手动改变晶格参数并画最低点附近能量。(B)用 \texttt{gulp1b.in} 的EAM势重复。(C)比较两者与实验晶格常数4.08 \AA{}，说明是否符合预期；评论两种模型总能量不同的含义。
\end{exercise}
\begin{solution}
原常规胞四原子的分数位置是 $(0,0,0),(0,.5,.5),(.5,0,.5),(.5,.5,0)$；手册LJ式 $v=A/r^{12}-B/r^6$ 给 $A=214180.2\,\text{eV}\,\text{\AA}^{12},B=625.482\,\text{eV}\,\text{\AA}^6$，截断40 \AA{}。opti加conp允许应力驱动晶格弛豫，原子弛豫与胞弛豫是两个变量组；single只评估指定构型。

为独立核算，程序对完整FCC邻居向量求 $e(a)=\tfrac12\sum_{\mathbf R\ne0,r<40\,\text{\AA}}v(r)$，按 \emph{每原子} 归一化。扫描与极小化得到 $a\simeq4.07802$ \AA{}，$e\simeq-3.93069$ eV/atom；常规胞总能量为 $4e$。采用未移位势与明示40 \AA{}截断，不声称逐软件选项复现GULP。其与4.08 \AA{}的接近只能检验这个量，不能证明势对缺陷、表面、弹性或所有相都正确。
\begin{center}\begin{tikzpicture}\begin{axis}[width=.7\linewidth,height=4.5cm,xlabel={$a$（\AA）},ylabel={$e$（eV/atom）}]
\addplot[main,thick] table[x=a_A,y=E_eV_per_atom,col sep=comma]{data/au-lj-static-lattice.csv};
\end{axis}\end{tikzpicture}\end{center}

EAM需给 $U=\sum_iF(\rho_i)+\tfrac12\sum_{i\ne j}\phi(r_{ij})$ 的F、密度核和成对项及截断、能量零点。取得原输入后，应保持相同FCC几何，扫描相同a区间，按相同胞归一化，再分别优化；不能只把LJ系数换个文件名。不同势的能量零点和原子参考态可不同，故两个绝对数值之差没有自动可比的物理意义。比较内聚能需明确共同的自由原子参考；比较同一势的构型能差则可消去常数。此处EAM数字未生成，真实缺失输入单列为执行限制。
\end{solution}

\originalsection{空位、弛豫与模型迁移}
\begin{exercise}\label{task:lab1-2}
（原Lab1题2）空位形成能 $E_f=E_{N-1,\mathrm{defect}}-(N-1)E_{N,\mathrm{perfect}}/N$。(A)用提供LJ势、不同超胞大小求未弛豫空位形成能，single而非opti；比较形成能与内聚能之比。(B)删去一个原子后用opti弛豫胞，比较形成能变化。(C)用原EAM势尽量准确复算，利用前两项的收敛经验并展示过程。(D)比较最终LJ结果与实验Au空位形成能0.9 eV，联系题1的晶格常数符合程度作评论。
\end{exercise}
\begin{solution}
完美超胞须先在同一势、截断与胞约定下求参考能。删原子后，形成能表达式相当于把该原子送回体相原子库，不能直接比较 $E_{N-1}-E_N$ 而漏掉原子数改变。周期边界意味着一个空位形成周期阵列；增大胞检查其相互作用，而非只增加完美体相原子数。

固定所有位置和胞时，纯成对势有一个透明的基准。设每原子体相能为e，空位位点与自身非零周期镜像的势和为 $S_{\rm img}$，删除该位点去掉与其他位点的键及自身镜像半计数键，故
\[
 E_f^{\rm unrelaxed}=-e+\tfrac12 S_{\rm img}.
\]
当胞无限大且势可积，镜像项趋零，$E_f\to-e$，形成能与内聚能\emph{绝对值}之比趋于1。新程序以相同40 \AA{}势计算 $n^3$ 个常规胞的阵列，$n=2,3,4,8$ 时 $E_f$ 分别约3.92181、3.92992、3.93056、3.93069 eV。它们保留镜像项，不是把“比为1”当成所有有限胞的精确结果。

原(B)要求允许原子和胞进一步弛豫。以未弛豫构型作可行初值，最小化能量不应增加；但若固定外压、固定体积或改变胞形，参考态和形成\emph{焓}约定必须相应一致。可检查空位近邻径向位移、最终最大力、残余应力及超胞收敛；仅有优化器退出码不足以判定完成。这里未运行GULP弛豫，不能虚构它降低了多少eV。

LJ未弛豫值约3.93 eV明显偏离0.9 eV，即使a几乎正确。它说明中心对势把断键能与体相内聚能强烈绑在一起，缺少环境密度的多体变化；原(C)的EAM能解除这一限制，但其准确性仍要实际检验原参数。空位弛豫也不保证任意势最终会吻合实验。
\end{solution}

\originalsection{表面能与两种尺寸收敛}
\begin{exercise}\label{task:lab1-3}
（原Lab1题3）(A)用原LJ势求FCC Au的(100)表面能，自选超胞和弛豫方法，写明公式并给eV/\AA$^2$；注意周期边界下的两个尺寸收敛问题并避免无谓增大体系。(B)用原EAM势重复；为节省时间可不作弛豫，仍须检查尺寸。
\end{exercise}
\begin{solution}
含两个等价表面的板片用 $\gamma=(E_{\rm slab}-Ne_{\rm bulk})/(2A_s)$，面积 $A_s$ 是单个表面面积。两项尺寸是\emph{固体板片厚度}与\emph{真空层宽度}；增加真空不能修复板片内部太薄的误差，增加板片也不能消去相邻镜像板片的作用。按(100)法向分层、保持平面方向尽小的周期重复即可减少原子数；板片与体相要用同一势、同一a与截断。

对未弛豫的(100)层间距 $h=a/2$、层面密度 $\rho_s=2/a^2$，令 $w_n$ 为参考层原子到第n层所有原子的成对势和。M层板片与无限体相的差给
\[
 \gamma_M=-\frac{\rho_s}{2}\sum_{n>0}\min(n,M)w_n.
\]
每种跨层键在切开界面时被断开的次数为n；薄板片对更远层的缺失只能按M计，因而出现最小值。独立代码在移除板片间作用的几何条件下求和，M为4、8、16、24时约 $.207407,.210275,.210754,.210760$ eV/\AA$^2$，末值约3.37674 J/m$^2$。这是未弛豫LJ表面，不是实测Au表面能，也未冒作EAM或已弛豫结果。

原(A)还应开放表层原子的位置，逐层记录弛豫深度，必要时比较固定平面a和允许应变的边界。有限厚度下两个表面可互相影响；对称施加上下表面处理能防止整体平移与不等价表面混入。原(B)保留同一板片和体相参考，只更换完整EAM输入，不能重用LJ的e作为EAM体相参考。
\end{solution}

\originalsection{势的适用性与成对最低距离}
\begin{exercise}\label{task:lab1-4}
（原Lab1题4）哪些其他问题适合LJ势，哪些适合EAM，哪些两者都不合适？
\end{exercise}
\begin{solution}
LJ保留短程排斥和色散吸引，可用于惰性气体或明确的简化非极性粒子模型及方法核验；真实分子还需要方向、键和电荷。EAM以环境密度改变嵌入能，常用于所拟合金属的缺陷、表面、弹塑性研究，但不是所有合金、相和高压状态的统一势。强方向共价键、显著电荷转移、化学反应或电子激发，通常不能靠这两种基本形式完整描述，应按要保留的自由度选择其他势或电子结构方法。拟合晶格成功不等于可迁移到所有这些问题。
\end{solution}
\begin{exercise}\label{task:lab1-5}
（原Lab1题5）LJ形式看似可用纸笔，为什么仍需计算机？求 $r_0(A,B)$，代入原Au参数，与题1晶格结果比较并解释；联系中心成对势常见的向外表面弛豫。
\end{exercise}
\begin{solution}
已核验的晶格和片段给 $r_0=(2A/B)^{1/6}=2.96892$ \AA{}；晶体最近邻距 $a/\sqrt2$ 约2.88360 \AA{}，不等于孤立对最低点，因为所有邻壳共同作用。非均匀构型还需同时优化许多原子和胞变量，手算一个二体极值无法处理这些收敛与几何约束。

在体相，较远吸引邻居可把最近邻距压到孤立对最低点以内；表面移去部分远邻后，剩余作用的平衡可能向较大局部距离移动，所以简单中心势常给向外表层变化。它不是任意表面、任何几何和全部参数的符号定理；切向/法向应分别分析，最终须计算层间距和力。这个机制解释为何原题要求计算与弛豫，而不是仅把 $\sqrt2r_0$ 当成所有表面或晶体的答案。
\end{solution}

\begin{exercise}\label{task:lab1-extra}
（原可选加分题）用EAM求FCC Au的 $C_{11},C_{12},C_{44}$；计算不同形变能量并拟合。
\end{exercise}
\begin{solution}
立方晶体每体积弹性能为 $\tfrac12 C_{11}(\epsilon_{xx}^2+\epsilon_{yy}^2+\epsilon_{zz}^2)+C_{12}(\epsilon_{xx}\epsilon_{yy}+\epsilon_{yy}\epsilon_{zz}+\epsilon_{zz}\epsilon_{xx})+2C_{44}(\epsilon_{xy}^2+\epsilon_{yz}^2+\epsilon_{zx}^2)$。单轴小应变、体积保持的正交应变、纯剪切分别提取独立组合；应同时取正负幅度、检查二次区间，并声明是否弛豫内部位置。剪切用工程应变 $\gamma_{xy}=2\epsilon_{xy}$ 时系数改变，不能漏因子4。原EAM输入未取得，数字保持未执行状态；方法和需保存的应变、能量、体积与力阈值完整保留。
\end{solution}
